\begin{satz}
{Eigenschaften der komplementären \textsl{N}(0,1)--Verteilungsfunktion}
{EigenschaftenDerKomplementaerenVerteilungsfunktionVonN01}
Seien $\gamma: \Omega \rightarrow \R$ $N(0,1)$--verteilt,
$u:\R_{\ge 0} \rightarrow [0,1], t \mapsto \Prim(\abs \gamma > t)$ und
schließlich
$\Phi: \R \rightarrow\!\  ]0,1[$ die Verteilungsfunktion von $N(0,1)$.\\
Dann gilt:
\begin{enumerate}[(1)]
\item \label{EigenschaftenDerKomplementaerenVerteilungsfunktionVonN011}
$u = 2(1 - \restrict{\Phi}{[0, \infty[}) = \sqrt{\frac 2 \pi}\int_\cdot ^\infty
\exp\left(-\frac{s^2}2 \right)ds$, $\Bild u =\!\  ]0,1]$, $u$ ist injektiv
und besitzt daher ein Inverses 
$u\inv:\!\ ]0,1] \rightarrow [0, \infty[$ mit $\Phi \circ u\inv 
= 1- \frac \cdot 2$.
\item \label{EigenschaftenDerKomplementaerenVerteilungsfunktionVonN012}
$u$ ist stetig differenzierbar auf $\R_{>0}$ mit 
\[\forall\,t \in \R_{>0}: u'(t) = 
-\sqrt{\frac 2 \pi}\exp\left(- \frac{t^2}2 \right)\text.\]
\item \label{EigenschaftenDerKomplementaerenVerteilungsfunktionVonN013}
Es gibt $(c_j)_{j \in \N_0}\in [0, \infty[^{\N_0}$ mit
\[\forall \, t \in\!\  ]0,1]: u\inv(t) = \sum_{j=0}^\infty
c_j(1-t)^{2j+1}\text.\]
\item \label{EigenschaftenDerKomplementaerenVerteilungsfunktionVonN014}
$\forall\, t \in \R_{>0}:
\frac{\sqrt{2\pi}}{(\pi-1)t + \sqrt{t^2+2\pi}}\exp\left(-\frac{t^2}2 \right)
\le u(t) \le \sqrt{\frac 2 \pi}\frac 1 t \exp\left(-\frac {t^2}2 \right)$.
\item \label{EigenschaftenDerKomplementaerenVerteilungsfunktionVonN015}
$\forall\,t\in [1,\infty[: u\inv\left(\frac 1 {t+1} \right)\ge \sqrt{\log(t)}$.
\item \label{EigenschaftenDerKomplementaerenVerteilungsfunktionVonN016}
$\forall\, t \in\!\  ]0, \frac 2 {13}]: 
 u \inv(t) \ge \sqrt{\log\left(1+\frac 1 t\right)}$.
\item \label{EigenschaftenDerKomplementaerenVerteilungsfunktionVonN017}
$\forall\, t \in \!\ ]0,\exp\left(- \frac 1 \pi \right)]:
u\inv(t) \le \sqrt{2\log\left(\frac 1 t\right)}$.
\end{enumerate}
\end{satz}
\begin{proof}
(1) Sei $t \in [0, \infty[$. Dann gilt:
\begin{align*}
u(t) &= 1 - \Prim(\abs \gamma \le t) = 1 - \Prim(-t \le \gamma \le t)
= 1 - (\Phi(t) - \Phi(-t))\\
&= 1 - \Phi(t) + 1- \Phi(t) = 2(1- \Phi(t)) = 2\left(\int_t^\infty \frac 1 {\sqrt{2 \pi}} \exp\left(- \frac{s^2}2 \right)ds\right)\\
& = \sqrt{\frac 2 \pi}\int_t ^\infty \exp\left(- \frac{s^2}2 \right)ds\text.
\end{align*}
Wegen $\Bild \Phi = \!\ ]0,1[$ gilt daher $0 < 2(1-\Phi(t)) = u(t)$ und 
wegen $t \ge 0$ und damit $\Phi(t) \ge \frac 1 2$ auch
$u(t) = 2(1-\Phi(t)) \le 1$.\\
Damit ist $\Bild u \subseteq \!\ ]0, 1]$.\\
Weiter ist $\Phi$ invertierbar und zu $y \in\!\  ]0,1]$ ist 
$1 - \frac y 2 \in [\frac 1 2, 1[\, \subseteq \Bild \Phi$, also:
\begin{align}
u\left(\Phi\inv\left(1- \frac y 2 \right) \right) = 2(1-(1-\frac y 2))=y\text.
\label{edkvvneq:1}
\end{align}
Damit ist $\Bild u =\!\  ]0,1]$.\\
Da $\Phi$ streng monoton wachsend ist, ist $u$ streng monoton fallend und somit
insbesondere injektiv, besitzt also ein Inverses auf $]0,1]$\\
Wegen $\eqref{edkvvneq:1}$ ist für alle $z\in\!\ ]0,1]$: 
$u(\Phi\inv(1-\frac z2))=z$, also $u\inv(z) = \Phi\inv(1- \frac z 2)$.\\
(2) Es ist $\Phi$ stetig differenzierbar auf $\R$ mit
\[\forall\,t\in\R:\Phi'(t)=\frac 1{\sqrt{2\pi}}\exp\left(-\frac{t^2}2
 \right)\text.\]
Daher gilt für alle $t \in \R_{>0}:$
\[u'(t)\overset{(1)}= 2(-\Phi'(t))=-\frac 2{\sqrt{2\pi}}\exp\left(-\frac{t^2}2
\right)=-\sqrt{\frac 2 \pi}\exp\left(-\frac{t^2}2 \right)\text.\]
(3) Sei $\erf: \R \rightarrow \R,
x \mapsto \frac 2 {\sqrt \pi}\int_0^x\exp(-t^2)dt$.\\
Dann gilt für alle $x \in \R$:
\begin{align*}
&\, \frac 1 2 + \frac 1 2 \erf\left(\frac x {\sqrt 2} \right)
= \frac 1 2 + \frac 1 {\sqrt \pi}\int_0^{\frac x {\sqrt 2}}\exp(-t^2)dt\\
\overset{\substack{\text{Subst.}\\ s \mapsto \frac s{\sqrt 2}}}=&\, \frac 1 2
+\frac 1{\sqrt\pi}\int_0^x \exp\left(- \frac {t^2}2 \right)\frac 1{\sqrt 2}dt\\
=&\, \spi \int_{-\infty}^0 \exp\left(-\frac {t^2}2 \right)dt 
+ \spi\int_0^x\exp\left(- \frac {t^2}2 \right)dt = \Phi(x)\text.
\end{align*}
Also $\Phi = \frac 1 2\left(1+\erf\left(\frac \cdot{\sqrt 2} \right)\right)$,
d. h. 
\begin{align}
\erf = 2\Phi(\sqrt 2 \cdot) - 1\text. \label{edkvvneq:2}
\end{align}
Für alle $\alpha\in\!\ ]-1,1[$ ist $\beta\coloneqq \frac 12(\alpha+1)\in\!\ ]0,1[$,
also gibt es ein $t \in \R$ mit $\Phi(t) = \beta$ und damit
\[\erf\left(\frac 1 {\sqrt 2}t \right) \overset{\eqref{edkvvneq:2}}= 2\Phi(t)-1
= 2\left(\frac 1 2 (\alpha +1)\right) -1 = \alpha \text.\]
Damit ist $\Bild \erf = \!\ ]-1,1[$.\\
Offensichtlich ist $\erf$ analytisch und es gilt $\erf' = \frac 2{\sqrt \pi}
\exp\left(- (\cdot^2) \right) > 0$.\\
Wegen der Stetigkeit des Integranden ist $\erf$ also streng monoton wachsend
und damit injektiv. Ferner ist damit $\erf\inv:\!\ ]-1,1[ \rightarrow \R$
ebenfalls analytisch.\\
Nach Setzung ist $\erf$ ungerade und damit nach
\refclaim{GeradeUndUngeradeFunktionen}{GeradeUndUngeradeFunktionen4} auch
$\erf\inv$ ungerade.\\
Sei $k \in \N$. Nach
\refclaim{GeradeUndUngeradeFunktionen}{GeradeUndUngeradeFunktionen2} und
\refclaim{GeradeUndUngeradeFunktionen}{GeradeUndUngeradeFunktionen3} ist dann
$(\erf\inv)^{(2k)}$ ungerade und damit nach
\refclaim{GeradeUndUngeradeFunktionen}{GeradeUndUngeradeFunktionen1}:
\begin{align}(\erf\inv)^{(2k)}(0) = 0 \text. \label{edkvvneq:3}\end{align}
Weiter gilt wegen $\id_{]-1,1[} = \erf \circ \erf\inv$ auch
$1 = (\erf\inv)' \cdot \erf' \circ \erf\inv$, also:
\begin{align}
\forall\, p \in \!\ ]-1,1[: (\erf\inv)'(p) = \frac 1 {\erf' ( \erf \inv(p))} 
=\frac {\sqrt \pi}2 \exp(\erf\inv(p)^2)\text. \label{edkvvneq:4}
\end{align}
Damit gilt aber auch für alle $p \in\!\  ]-1,1[$:
\begin{align*}(\erf\inv)''(p) &\overset{\eqref{edkvvneq:4}}=
2\cdot\frac{\sqrt \pi}2 \erf\inv(p)(\erf\inv)'(p)
\exp\left(\erf\inv(p)^2\right)\\
&\overset{\eqref{edkvvneq:4}} = 
2\cdot\erf\inv(p)\cdot ((\erf\inv)'(p))^2\text.\end{align*}
Ferner gilt $(\erf\inv)'(0)= \frac {\sqrt\pi}2 \exp(\erf\inv(0)^2) > 0$, also
gilt nach~\refclaim{GeradeUndUngeradeFunktionen}{GeradeUndUngeradeFunktionen5}:
\begin{align}
(\erf\inv)^{(2k+1)}(0) \ge 0 \text.\label{edkvvneq:5}
\end{align}
Nach~\cite{Dom} (dort Abschnitt 4) ist die Taylorreihe von $\erf\inv$ mit
Entwicklungspunkt $0$ auf $]-1,1[$ konvergent. Wegen der Analytizität von
$\erf\inv$ stellt die besagte Reihe also auf $]-1,1[$ bereits $\erf\inv$ dar.\\
Also liefert die Taylorentwicklung um $0$ für alle $p \in \!\ ]-1,1[$:
\begin{align}
\erf\inv(p) = \sum_{l=0}^\infty \frac{(\erf\inv)^{(l)}(0)}{l!}p^l
\overset{\eqref{edkvvneq:3}}=\sum_{l=0}^\infty
\frac{(\erf\inv)^{(2l+1)}(0)}{(2l+1)!}p^{2l+1}\text.\label{edkvvneq:6}
\end{align}
Wegen~\eqref{edkvvneq:2} gilt weiter: $\Phi\inv=\sqrt 2\erf\inv(2\cdot -1)$.\\
Damit für alle $q \in \!\ ]0,1]$ wegen $1-\frac q 2 \in\!\  ]-1,1[$:
\begin{align*}
u\inv(q) &\overset{(1)}= \Phi\inv\left(1- \frac q 2\right)
= \sqrt 2 \erf\inv\left(2\left(1-\frac q 2 \right)-1 \right)
= \sqrt 2 \erf\inv(1-q)\\
&\overset{\eqref{edkvvneq:6}}= \sum_{l=0}^\infty
 \sqrt 2 \frac{(\erf\inv)^{(2l+1)}(0)}{(2l+1)!}(1-q)^{2l+1}\text.
\end{align*}
Nach~\eqref{edkvvneq:5} sind alle Koeffizienten dieser Reihendarstellung
nichtnegativ, also ist (3) gezeigt.\\
(4) Seien \[V_0 : [0, \infty[ \rightarrow \R, 
x \mapsto 2\exp\left(x^2\right)\int_x^\infty \exp\left(-s^2\right)ds \]
und 
\[g_\pi:[0,\infty[\rightarrow\R,x\mapsto \frac{\pi}{(\pi-1)x+\sqrt{x^2+\pi}}
\text.\]
Nach~\cite{RW} (Notation wie dort) gilt dann:
\begin{align}
g_\pi\le V_0\text.
\label{edkvvneq:7}
\end{align}
Sei $t \in \R_{>0}$. Dann gilt:
\begin{align*}
& \sqrt{\frac 2 \pi}\frac{\pi}{(\pi-1)t + \sqrt{t^2+2\pi}} =
\sqrt{\frac 2 \pi}\cdot \frac{\pi}{(\pi-1)\frac{\sqrt 2}{\sqrt 2}t
+ \sqrt{2\frac{t^2}{2} + 2\pi}}\\
 = & \,
\sqrt{\frac 2 \pi} \cdot \frac \pi {\sqrt 2\left((\pi-1)\frac 1{\sqrt 2}t  +
\sqrt{\left(\frac 1 {\sqrt{2}}t\right)^2 + \pi} \right)} 
= \sqrt{\frac 2 \pi} \frac 1 {\sqrt 2} g_\pi\left(\frac 1 {\sqrt 2} t \right)\\
\overset{\eqref{edkvvneq:7}}\le & \, \sqrt{\frac 2 \pi} \frac 1 {\sqrt 2} 
V_0\left(\frac 1 {\sqrt 2} t \right) = 
\sqrt{\frac 2 \pi}\exp\left(\frac {t^2} 2 \right)
\int_{\frac 1 {\sqrt 2} t}^\infty \exp\left(-s^2 \right)\sqrt 2 ds\\
= & \,\sqrt{\frac 2 \pi}\exp\left(\frac {t^2} 2 \right)
\int_{\frac 1{\sqrt 2}t}^\infty 
\exp\left(-\frac{(\sqrt 2s)^2}{2}\right)\sqrt 2 ds\\
\overset{\substack{\text{Subst.}\\ s \mapsto \sqrt 2 s}}=&\,
\sqrt{\frac 2 \pi}\exp\left(\frac {t^2} 2 \right)
\int_{t}^\infty \exp\left(- \frac{s^2}2 \right)ds
\overset{(1)}= \exp\left( \frac{t^2} 2 \right) u(t)\text,\end{align*}
also nach Umstellung und Kürzung:
\[\frac{\sqrt{2\pi}}{(\pi-1)t+\sqrt{t^2+2\pi}}\exp\left(-\frac{t^2}2\right)
 \le u(t)\text.\]
Weiter gilt auch:
\begin{align*}
t\cdot u(t) &\overset{(1)}= 
\sqrt{\frac 2 \pi} \int_t^\infty t \exp\left(- \frac {s^2}2 \right)ds
\le \sqrt{\frac 2 \pi} \int_t^\infty s \exp\left(- \frac {s^2}2 \right)ds\\
&= \sqrt{\frac 2 \pi}\exp\left(- \frac{t^2}2 \right)\text,
\end{align*}
also
\[u(t)\le\sqrt{\frac 2 \pi} \frac 1 t \exp\left(- \frac{t^2}2 \right)\text.\]
(5) Annahme: es gibt ein $t \in [1, \infty[$ mit 
$u\inv\left(\frac 1 {t+1} \right) < \sqrt{\log(t)}$.\\
Dann gilt:
\begin{align*}
&\frac 1 {t+1} = u\left(u\inv\left(\frac 1{t+1} \right)\right)\\
 \overset{(4)}\ge &\,
\frac{\sqrt{2\pi}}{(\pi-1)u\inv\left(\frac 1{t+1}\right)
+ \sqrt{\left(u\inv\left(\frac 1{t+1}\right)\right)^2+2\pi}}
\exp\left(-\frac 1 2 \left(u\inv\left(\frac 1{t+1}\right)\right)^2\right)\\
 >&\, \frac{\sqrt{2\pi}}{(\pi-1)\sqrt{\log(t)} 
+ \sqrt{\log(t)+2\pi}}\exp\left(-\frac{\log(t)}2\right)\\
=&\, \frac{\sqrt{2\pi}}{(\pi-1)\sqrt{\log(t)} 
+ \sqrt{\log(t)+2\pi}}\frac 1 {\sqrt t}\text.
\end{align*}
Somit auch
\[ \frac{\sqrt t}{t+1} \left((\pi-1)\sqrt{\log(t)} 
+ \sqrt{\log(t)+2\pi} \right) > \sqrt{2\pi}\text.\]
Sei $k: [1, \infty[ \rightarrow \R, 
s \mapsto \frac{\sqrt s}{s+1} \left((\pi-1)\sqrt{\log(s)} 
+ \sqrt{\log(s)+2\pi} \right)$.\\
Dann gilt $k(t) > \sqrt{2\pi}$ und nach
\refclaim{Numeriklemma}{Numeriklemma1} auch $k(t) <\sqrt{2\pi}$. $\lightning$\\
(6) Kontraposition:
Sei $t\in \!\ ]0, 1]$ mit $u\inv(t)<\sqrt{\log\left(1+\frac 1 t\right)}$.\\
Dann gilt:
\begin{align*}
t & = u\left(u\inv(t)\right) \overset{(4)}\ge
\frac{\sqrt{2\pi}}{(\pi-1)u\inv(t) 
+ \sqrt{u\inv(t)^2+2\pi}}\exp\left(-\frac{u\inv(t)^2}2\right)\\
& >  \frac{\sqrt{2\pi}}{(\pi-1) \sqrt{\log\left(1 +\frac 1 t\right)}
+ \sqrt{\log\left(1 +\frac 1 t\right)+2\pi}}
\exp\left(-\frac{\log\left(1 +\frac 1 t\right)}2\right)\\
& = \frac{\sqrt{2\pi}}{(\pi-1) \sqrt{\log\left(1 +\frac 1 t\right)}
+ \sqrt{\log\left(1 +\frac 1 t\right)+2\pi}}
\frac 1 {\sqrt{1+\frac 1 t}}\text.
\end{align*}
Also
\[t\sqrt{1+\frac 1 t} \left((\pi-1) \sqrt{\log\left(1 +\frac 1 t\right)}
+ \sqrt{\log\left(1 +\frac 1 t\right)+2\pi}\right)  > \sqrt{2\pi}\text.\]
Sei $h: [1, \infty[\!\  \rightarrow \R, s \mapsto
\frac{\sqrt{s+1}}s
\bigg( (\pi-1)\sqrt{\log(1+s)} + \sqrt{\log(1+s) + 2\pi} \bigg)$.\\
Dann ist also $\frac 1 t \in [1, \infty[$ und 
$h\left(\frac 1 t \right) > \sqrt{2\pi}$.\\
Nach~\refclaim{Numeriklemma}{Numeriklemma2} gilt daher $\frac 1 t <\frac{13}2$,
also $t > \frac 2 {13}$, d.h. $t \in\!\  ]\frac 2 {13}, 1]$ und damit
$t \notin \!\ ]0, \frac 2 {13}[$.\\
(7) Sei $t\in\!\  ]0,1]$ mit $u\inv(t) > \sqrt{2\log\left(\frac 1 t\right)}$.\\
Dann gilt:
\begin{align*}
t &= u\left(u\inv(t)\right) \overset{(4)}\le 
\sqrt{\frac 2 \pi} \frac 1 {u\inv(t)}
\exp\left(- \frac 1 2 \left(u\inv(t) \right)^2 \right)\\
&< \sqrt{\frac 2 \pi} \frac 1 {\sqrt{2\log\left(\frac 1 t\right)}}
\exp\left(- \frac 1 2 2\log\left(\frac 1 t\right) \right)
= \sqrt{\frac 2 \pi} \frac 1 {\sqrt{2\log\left(\frac 1 t\right)}} t\text.
\end{align*}
Daher also
$\sqrt{2\log\left(\frac 1 t \right)} < \sqrt {\frac 2 \pi}$, also 
 $\log\left(\frac 1 t \right) < \frac 1 \pi$
und somit $t > \exp\left(- \frac 1 \pi \right)$.
\end{proof}